% ---
% Reemplaza líneas 2234–2575 de Docs/Technical_Report.tex.
% Decisiones que lo respaldan: D-006, D-008, D-009, D-010, D-013, D-014, D-016, D-017,
% D-018, D-020, D-021, D-022, D-023, D-024, D-025, D-026, D-027, D-028, D-030, D-031,
% D-032, D-033, D-035, D-036, D-037, D-039, D-040, H-019, H-020, H-022, H-025, H-032,
% H-033, H-038, H-052.
% ---

\subsection{Requerimientos del sistema}
\label{sec:requerimientos}

El sistema desarrollado en el presente trabajo constituye un prototipo de 
investigación cuyo diseño responde a restricciones tanto técnicas como 
experimentales. Los requerimientos aquí presentados fueron definidos en 
concordancia con los objetivos específicos del proyecto y con las 
características del dataset CBIS-DDSM, y se organizan en requerimientos 
funcionales, que describen las capacidades que el sistema debe proveer,
y requerimientos no funcionales, que establecen las condiciones de calidad, 
reproducibilidad y viabilidad bajo las cuales el sistema debe operar.

% -----------------------------------------------------------
% 4.2.1 Requerimientos funcionales
% -----------------------------------------------------------
\subsubsection{Requerimientos funcionales}

Se presentan los requerimientos funcionales de mayor relevancia
para justificar las decisiones de diseño del pipeline,
organizados por las etapas que determinan la arquitectura
del sistema. Cada requerimiento está identificado con un
código RF-XX y vinculado al objetivo específico que lo origina.

\begingroup
\renewcommand{\arraystretch}{1.35}
\begin{longtable}{@{}lp{9.8cm}cc@{}}

\caption{Requerimientos funcionales del sistema}
\label{tab:req_funcionales} \\

% ── Encabezado de la PRIMERA página
\toprule
\textbf{ID} & \textbf{Descripción} &
\textbf{OE} & \textbf{Prior.} \\
\midrule
\endfirsthead

% ── Encabezado de páginas SIGUIENTES
\multicolumn{4}{c}{%
  \small\textit{Tabla \thetable{} -- continuación}} \\[4pt]
\toprule
\textbf{ID} & \textbf{Descripción} &
\textbf{OE} & \textbf{Prior.} \\
\midrule
\endhead

% ── Pie de páginas intermedias
\midrule
\multicolumn{4}{r}{%
  \small\textit{Continúa en la siguiente página}} \\
\endfoot

% ── Pie de la ÚLTIMA página
\bottomrule
\endlastfoot

\multicolumn{4}{@{}l}{%
  \cellcolor{gray!10}\small\textit{%
  Módulo 1 -- Lectura de datos}} \\[2pt]

RF-01 &
Leer archivos DICOM del dataset CBIS-DDSM mediante
\texttt{pydicom}, extrayendo el arreglo de píxeles en
16 bits y los metadatos clínicos.
& OE-1 & Alta \\

RF-02 &
Unificar los cuatro archivos CSV del dataset en un
DataFrame consolidado con etiqueta binaria
(BENIGN\,=\,0 / MALIGNANT\,=\,1).
& OE-1 & Alta \\

RF-03 &
Tratar de forma independiente masas y calcificaciones
como subproblemas de clasificación independientes.
& OE-1, OE-3 & Alta \\

\multicolumn{4}{@{}l}{%
  \cellcolor{gray!10}\small\textit{%
  Módulo 2 -- Preprocesamiento}} \\[2pt]

RF-04 &
Alinear cada máscara ROI con su mamografía según el tipo
de desajuste (remuestreo por vecino más cercano o anclaje
sin escalar) y reducirla a su componente conexa mayor.
& OE-1 & Alta \\

RF-05 &
Generar una variante con CLAHE sobre la mamografía completa
convertida a 8 bits, solo para cuantificar su efecto sobre las
características; la ruta de extracción usa la imagen sin CLAHE.
& OE-1 & Alta \\

RF-06 &
Recortar la ROI con la caja de la máscara ampliada en 20
píxeles por lado, a resolución original y con las
intensidades originales de 16 bits.
& OE-1 & Alta \\

\multicolumn{4}{@{}l}{%
  \cellcolor{gray!10}\small\textit{%
  Módulo 3 -- Extracción de características}} \\[2pt]

RF-07 &
Extraer 67 características radiómicas con PyRadiomics
(primer orden, forma 2D, GLCM y GLRLM) sobre la imagen
original, con \texttt{binCount} = 32.
& OE-1 & Alta \\

\multicolumn{4}{@{}l}{%
  \cellcolor{gray!10}\small\textit{%
  Módulo 4 -- Selección y escalado}} \\[2pt]

RF-08 &
Seleccionar $k = 12$ características por F-test ANOVA,
aceptando una solo si su correlación absoluta con las ya
elegidas es $\le 0.95$, con el selector ajustado
únicamente sobre el conjunto de entrenamiento.
& OE-1, OE-2 & Alta \\

RF-09 &
Transformar cada característica en un ángulo
$\pi\,\hat F(r) \in [0,\pi]$, con $\hat F$ su distribución empírica, mediante una
transformación por cuantiles ajustada sobre el conjunto
de entrenamiento.
& OE-1, OE-2 & Alta \\

\multicolumn{4}{@{}l}{%
  \cellcolor{gray!10}\small\textit{%
  Módulo 5 -- Transformación (5 condiciones)}} \\[2pt]

RF-10 &
\textit{Baseline} crudo: vector seleccionado y escalado
directo al MLP sin transformación adicional.
& OE-2 & Alta \\

RF-11 &
\textit{Baseline} lineal: PCA sobre el vector
seleccionado y escalado, conservando $k$ componentes.
Actúa como control nulo: una PCA de $k$ a $k$ dimensiones
es una rotación.
& OE-2 & Alta \\

RF-12 &
\textit{Baseline} no lineal: Kernel PCA con kernel RBF
sobre el vector seleccionado, con el ancho de banda
$\gamma$ fijado por la heurística de la mediana sobre
el entrenamiento.
& OE-2 & Alta \\

RF-13 &
Condición C4: \texttt{zz\_feature\_map} de segundo orden
seguido de \texttt{real\_amplitudes} con parámetros
$\boldsymbol\theta$ fijos (semilla 42, sin entrenar),
devolviendo los valores esperados exactos
$\langle Z_i\rangle \in [-1,1]^k$.
& OE-2 & Alta \\

RF-14 &
Condición C5: \texttt{pauli\_feature\_map} con operadores
$(Y, ZZ)$, con el mismo \textit{ansatz} y los mismos
$\boldsymbol\theta$ que C4.
& OE-2 & Alta \\

\multicolumn{4}{@{}l}{%
  \cellcolor{gray!10}\small\textit{%
  Módulo 6 -- Análisis de separabilidad}} \\[2pt]

RF-15 &
Calcular el kernel cuántico de fidelidad, solo con el
\textit{feature map}, sobre 200 lesiones de entrenamiento por
subconjunto, y su alineamiento centrado con las etiquetas,
comparado contra una prueba de permutación.
& OE-3 & Alta \\

RF-16 &
Calcular la \textit{geometric difference} regularizada entre
el kernel cuántico y una batería de kernels clásicos, con el
protocolo de Huang et al., junto con la dimensión efectiva:
$g_{CQ}=\sqrt{\bigl\|\sqrt{K_Q}\,\sqrt{K_C}\,(K_C+\lambda I)^{-2}\sqrt{K_C}\,\sqrt{K_Q}\bigr\|_\infty}$.
& OE-3 & Alta \\

RF-17 &
Medir el índice de Davies-Bouldin, la razón de Fisher
multivariante y el alineamiento del kernel lineal sobre la
representación de cada condición, separado para masas y
calcificaciones.
& OE-3 & Alta \\

\multicolumn{4}{@{}l}{%
  \cellcolor{gray!10}\small\textit{%
  Módulo 7 -- Clasificación}} \\[2pt]

RF-18 &
Entrenar el mismo MLP ($k \to 32 \to 16 \to 1$) en las
5 condiciones, con idéntica arquitectura, hiperparámetros,
protocolo y semilla.
& OE-4 & Alta \\

RF-19 &
Entrenar con ponderación de clases, con pesos
$w_c = N/(C\,N_c)$ (ecuación~\eqref{eq:class_weights}) en la
entropía cruzada binaria para compensar el desbalance
observado.
& OE-4 & Alta \\

\multicolumn{4}{@{}l}{%
  \cellcolor{gray!10}\small\textit{%
  Módulo 8 -- Evaluación comparativa}} \\[2pt]

RF-20 &
Reportar AUC-ROC, F1-Score, Accuracy, Precision y
Recall por condición, separado para masas y
calcificaciones.
& OE-5 & Alta \\

RF-21 &
Analizar la correlación entre ranking de separabilidad
(OE-3) y ranking de clasificación (OE-5) de forma
descriptiva, con la distribución exacta de $\rho_s$
para $n = 5$.
& OE-5 & Alta \\

RF-22 &
Documentar restricciones del simulador: tiempo por
muestra, qubits máximos, impacto de \textit{shots} y
costo del \textit{parameter-shift}.
& OE-6 & Alta \\

\end{longtable}
\endgroup

% -----------------------------------------------------------
% 4.2.2 Requerimientos no funcionales
% -----------------------------------------------------------
\subsubsection{Requerimientos no funcionales}

Los requerimientos no funcionales establecen las condiciones
bajo las cuales el sistema debe operar para garantizar la
validez científica y reproducibilidad de los resultados.

\begingroup
\renewcommand{\arraystretch}{1.35}
\begin{longtable}{@{}llp{9.5cm}@{}}

\caption{Requerimientos no funcionales del sistema}
\label{tab:req_no_funcionales} \\

\toprule
\textbf{ID} & \textbf{Categoría} &
\textbf{Descripción} \\
\midrule
\endfirsthead

\multicolumn{3}{c}{%
  \small\textit{Tabla \thetable{} -- continuación}} \\[4pt]
\toprule
\textbf{ID} & \textbf{Categoría} &
\textbf{Descripción} \\
\midrule
\endhead

\midrule
\multicolumn{3}{r}{%
  \small\textit{Continúa en la siguiente página}} \\
\endfoot

\bottomrule
\endlastfoot

RNF-01 & Reproducibilidad &
Fijar semillas aleatorias en \texttt{numpy}, \texttt{PyTorch}
y \texttt{Qiskit} al inicio de cada experimento (semilla 42).
Las semillas de los \textit{folds} (125 en masas y 168 en
calcificaciones) se eligen por balance de clases. \\

RNF-02 & Reproducibilidad &
Declarar los parámetros de cada módulo en su \textit{notebook}
y guardar los parámetros ajustados junto a los datos procesados
(\texttt{x\_\textless tipo\textgreater.json} en M4, \texttt{emb\_meta.json} en M5). \\

% respaldo: «The data were split into a training set and a testing set based on the BI-RADS category.» (Lee2017CBISDDSM, p. 5)
RNF-03 & Reproducibilidad &
Usar la partición oficial train/test sin reordenarla, que los
autores del CBIS-DDSM estratificaron por categoría BI-RADS
\cite{Lee2017CBISDDSM}, con la partición real de 2\,864 / 704 lesiones (80.3\,\% / 19.7\,\%);
el test se evalúa una sola vez por modelo y la validación
cruzada usa 5 \textit{folds} agrupados por paciente dentro del
train. \\

RNF-04 & Reproducibilidad &
Ajustar el selector, el escalador por cuantiles y la
estandarización del MLP solo con datos de entrenamiento para
transformar el conjunto de prueba, previniendo \textit{data leakage}. \\

RNF-05 & Rendimiento &
Calcular los valores esperados de forma exacta para el
experimento, y cuantificar aparte el efecto de los
\textit{shots} muestreando $n$ disparos de las probabilidades
exactas (de 64 a 65\,536, incluido 1\,024). \\

RNF-06 & Rendimiento &
Calcular la matriz de kernel cuántico sobre submuestra de
200 lesiones de entrenamiento por subconjunto para viabilidad
computacional del cálculo $O(N^2)$ de \textit{geometric difference}
y KTA. \\

RNF-07 & Escalabilidad &
Circuito parametrizable en $k$, \textit{feature map} y
\texttt{reps}. Punto de operación $k = 12$ y \texttt{reps} = 1,
con $k = 8$ y \texttt{reps} = 2 como análisis de sensibilidad. \\

RNF-08 & Mantenibilidad &
Código organizado en ocho módulos M1--M8, un \textit{notebook}
por módulo. \\

RNF-09 & Propiedad intelectual &
Registro del trabajo en INDAUTOR. \\

\end{longtable}
\endgroup

\begin{table}[H]
\centering
\caption{Correspondencia entre fases CRISP-DM y módulos
         del sistema}
\label{tab:crisp_modulos}
\renewcommand{\arraystretch}{1.3}
\begin{tabularx}{\textwidth}{@{}llX@{}}
\toprule
\textbf{Fase CRISP-DM} & \textbf{Módulo} &
\textbf{Actividad principal} \\
\midrule
1. Comprensión del negocio  &
  ---   &
  Definición del problema de prediagnóstico,
  formulación de la pregunta de investigación
  y criterios de separabilidad. \\
2. Comprensión de los datos &
  M1    &
  Lectura DICOM, análisis exploratorio del
  dataset CBIS-DDSM (EDA). \\
3. Preparación de los datos &
  M2, M3, M4 &
  Preprocesamiento de imagen, extracción de
  features y selección/escalado. \\
4. Modelado &
  M5, M7 &
  Implementación de las 5 condiciones de
  transformación y entrenamiento del MLP. \\
5. Evaluación &
  M6, M8 &
  Análisis de separabilidad, métricas de
  clasificación, comparación clásico frente a
  cuántico y restricciones del simulador. \\
6. Despliegue &
  ---   &
  Documentación de hallazgos, reporte técnico
  y registro INDAUTOR. \\
\bottomrule
\end{tabularx}
\end{table}
